This paper evaluated a DQN-VQC agent against a parameter-matched MLP-DQN
baseline on CartPole-v1. We summarize findings against four research questions.

\textbf{RQ1.} Under a 22--23 parameter budget and 100 training episodes,
DQN-MLP and DQN-VQC are indistinguishable under 5 seeds (both $\approx 9.4$,
below random at 20.65). Extended 500-episode training recovered learning in 1
of 3 MLP seeds (greedy 69.45) while VQC remained below random.

\textbf{RQ2.} Double and Dueling DQN did not improve outcomes at 100 episodes;
all variants collapsed to the same deterministic policy. Training budget, not
the value-estimation technique, dominates results at this scale.

\textbf{RQ3.} Batched QNode execution via PennyLane parameter
broadcasting~\cite{bergholm2018pennylane} reduced VQC simulation overhead from
$44\times$ to $5\times$ relative to MLP (both on CPU). The multiplier conflates
PyTorch-vs-simulator overheads.

\textbf{RQ4.} Under matched supervised protocols, MLP L1 represents the teacher
from raw features (500.00, 5 seeds) while VQC L1 does not (96.53). VQC L3 raw
reaches 321.87, showing depth helps. With fuzzy membership features, VQC L1 and
L3 achieve 500.00. Critically, VQC L1 DQN with the same features collapses to
9.60 (5 seeds), indistinguishable from raw-feature DQN. Since the MLP also fails
under DQN despite solving raw-feature distillation (500.00), the bottleneck
is not VQC-specific: it is a general inability of the DQN reward signal — at
this budget, with this epsilon schedule, on this environment — to extract
policies that the models demonstrably can represent under supervised
optimization.

Future work should extend to additional environments, increase seeds further,
and investigate hybrid strategies combining fuzzy-guided exploration with
curriculum learning for small VQC models.
